
% =====================================================================
\section{Which fermion pairs exist}
\label{sec:pairs}
% =====================================================================

A superfluid of fermions is a condensate of pairs. Before any dynamics is introduced, the matrices of
Sec.~\ref{sec:matrices} decide which pairs of two quanta of the field can exist at all. This section
derives that list. Its result, the pair along the hidden axis, is the object that the rest of the
paper studies.

\subsection{No Lorentz-invariant pair}

A pair operator built from two field operators of the same kind has the form
$\psi^TX\psi=\sum_{a,b}\psi_aX_{ab}\psi_b$ with a $16\times16$ matrix $X$. Because the components
anticommute, only the antisymmetric part of $X$ contributes; a symmetric $X$ gives an operator that
vanishes identically. The matrices
\begin{equation}
X=C\,\gamma^{A_1}\cdots\gamma^{A_r}\qquad(A_1<\dots<A_r),
\end{equation}
with $C=\Cm$ or $C=\Cp$, transform under Lorentz transformations as the components of antisymmetric
tensors of rank $r$, because by Eq.~(\ref{eq:Ctranspose})
$S^TC\,\gamma^{A_1}\cdots\gamma^{A_r}S=C\,S^{-1}\gamma^{A_1}\cdots\gamma^{A_r}S$. The program
determines the symmetry of all 256 of them for both matrices $C$ (check
\SIG{\texttt{44-K1}}{\texttt{71-K1}}). Table~\ref{tab:pairsym} gives the result: all matrices of a
given rank have the same symmetry, and 120 of the 256 are antisymmetric, as many as there are
independent products $\psi_a\psi_b$ with $a<b$.

\begin{table}
\caption{\label{tab:pairsym}Symmetry of the matrices $C\,\gamma^{A_1}\cdots\gamma^{A_r}$ of rank $r$
for the two transposition matrices of the \SIG{$4+4$}{$7+1$} spacetime. Only an antisymmetric matrix
gives a pair operator $\psi^TX\psi$ that does not vanish. The pattern is the same for the companion
spacetime.}
\begin{ruledtabular}
\begin{tabular}{cccc}
rank $r$ & matrices & $\Cm_{\tg}\gamma\cdots\gamma$ & $\Cp_{\tg}\gamma\cdots\gamma$ \\
\hline
{{ROWS:pair_symmetry_table|cols=\SIG{0,1,2,3}{0,1,4,5}}}
\end{tabular}
\end{ruledtabular}
\end{table}

\emph{Theorem \SIG{3}{4}.} The Lorentz-scalar pair operators $\psi^T\Cm\psi$ (rank 0) and
$\psi^T\Cm\omega\,\psi=\pm\psi^T\Cp\psi$ (rank 8) vanish identically (check
\SIG{\texttt{44-K2}}{\texttt{71-K2}}). A single 16-component field has no Lorentz-invariant fermion
pair. The pair operators that exist are, in terms of $\Cm$, a vector ($r=1$, eight components), a
two-form (28), a five-form (56) and a six-form (28).

The theorem is a statement about the eight-dimensional Lorentz group. It does not exclude a pair that
is a scalar for an observer who sees only four of the eight dimensions, and that is the possibility
pursued below.

\subsection{The pairs of two particles at rest}
\label{sec:pairsrest}

A superfluid forms in a medium, and the medium singles out a rest frame. The relevant question is
then: which pairs of two identical particles at rest, in a state of zero orbital angular momentum (an
\emph{$s$ wave}), are allowed, and how do they transform under the rotations that the medium leaves
intact? One particle at rest has eight internal states. The Pauli principle admits the
$8\cdot7/2=28$ antisymmetric combinations, which are described by the antisymmetric $8\times8$
\emph{pair matrices}.

%%<1644>
We work in the quantized sector, with fields independent of $x5,x6,x7$. A particle at rest satisfies
$\gamma^4u=iu$. The symmetry that the medium leaves intact is the rotation group $SO(4)$ of the four
space coordinates $x0,\dots,x3$, the rotation group $SO(3)_T$ of the three additional times, and the
sign of the grading $W$. The program classifies the 28 pairs (checks \texttt{44-Q1} to
\texttt{44-Q4}).
\begin{itemize}
\item The pairs that are singlets under all rotations of the four space dimensions are six in number,
two triplets of $SO(3)_T$. None of them is a singlet of $SO(3)_T$ as well.
\item Ten pairs have ordinary spin 0, that is, are invariant under the rotations of $x1,x2,x3$.
\item \emph{Exactly one} pair is a singlet both of the ordinary rotations and of the rotations of
the additional times. It is the two-particle content of the operator $\psi^T\Cm\gamma^0\psi$, the
component of the vector pair along the hidden axis.
\end{itemize}
The rotation group of the four space coordinates contains the ordinary rotations and mixes the hidden
axis with ordinary space; it is not an internal symmetry of ordinary physics. It must not be confused
with the internal $SO(4)$ of Sec.~\ref{sec:internal}, which appears only for fields that do not
depend on $x0$.
%%</1644>
%%<1671>
On particles at rest the rotation generators act as $\tfrac12\tau^i\tau^j$, and the 28 antisymmetric
$8\times8$ matrices are spanned by the seven $\tau^j$ and the twenty-one products $\tau^i\tau^j$ with
$i<j$, where the indices run over the seven space coordinates. The program classifies the pairs
(checks \texttt{71-Q1} to \texttt{71-Q10}).
\begin{itemize}
\item The only matrix that commutes with all rotations is the unit matrix, which is symmetric. There
is no $s$-wave pair that is a singlet under the rotations $SO(7)$ of the seven space dimensions.
\item Under $SO(7)$ the 28 pairs form a vector, spanned by the seven $\tau^j$, and an adjoint,
spanned by the twenty-one $\tau^i\tau^j$.
\item Split the seven space coordinates into the ordinary ones, $x1,x2,x3$, whose rotations $SO(3)$
generate the ordinary spin, and the four extra axes $x0,x5,x6,x7$, whose rotations form the group
$SO(4)$, locally $SU(2)_A\times SU(2)_B$. Exactly ten pairs have ordinary spin 0: the four matrices
$\tau^0,\tau^5,\tau^6,\tau^7$, which form a four-vector of $SO(4)$, the representation $(2,2)$ of
$SU(2)_A\times SU(2)_B$; and the six matrices $\tau^a\tau^b$ with $a<b$ in $\{0,5,6,7\}$, which form
$(3,1)\oplus(1,3)$.
\item \emph{Exactly one} pair is a singlet both of the ordinary rotations and of the rotations of
$x5,x6,x7$: the pair matrix $\tau^0$.
\item The corresponding field operators are $\psi^T\Cm\gamma^j\psi$ for the vector and
$\psi^T\Cp\gamma^i\gamma^j\psi$ for the adjoint.
\end{itemize}
A single particle at rest transforms as $(2;2,1)\oplus(2;1,2)$ under
$SO(3)\times SU(2)_A\times SU(2)_B$: it has ordinary spin $\tfrac12$ and is a doublet of one of the
two groups of the extra space.
%%</1671>

\subsection{The pair along the hidden axis}
\label{sec:hiddenpair}

In both spacetimes there is a \emph{vector pair operator},
\begin{equation}
\Delta_{\tg}^{A}\equiv\psi_{\tg}^{T}\,\Cm_{\tg}\,\gamma_{\tg}^{A}\,\psi_{\tg},\qquad A=0,\dots,7,
\label{eq:vectorpair}
\end{equation}
and its component along the hidden coordinate $x0$,
\begin{equation}
\Delta_{\tg}^{0}=\psi_{\tg}^{T}\,\Cm_{\tg}\,\gamma_{\tg}^{0}\,\psi_{\tg},
\label{eq:hidden}
\end{equation}
has three properties. (i) Its two-particle content has ordinary spin 0. (ii) It is the only $s$-wave
pair that is invariant under the ordinary rotations and under the rotations of $x5,x6,x7$
(Sec.~\ref{sec:pairsrest}). (iii) It is invariant under all 21 generators $\Sigma^{AB}$ with
$A,B\in\{1,\dots,7\}$, and only the 7 generators $\Sigma^{0B}$ that involve the hidden axis change it
(check \SIG{\texttt{44-K3}}{\texttt{71-K3}}). Property (iii) follows from Eq.~(\ref{eq:vector}): a
transformation generated by $\Sigma^{AB}$ changes $\Delta^{0}$ by
$-\psi^T\Cm[\Sigma^{AB},\gamma^0]\psi$, which vanishes unless $A$ or $B$ is 0.

In other words, $\Delta^{0}$ is a \emph{scalar} under the Lorentz group of the seven coordinates other
than $x0$, and in particular under the Lorentz group of ordinary spacetime $(x1,x2,x3;x4)$, although
it is a component of a vector of the eight-dimensional spacetime. This is the property for which the
hidden axis is needed. A scalar pair of the full eight-dimensional group does not exist; but a vector
along a dimension that ordinary physics does not see is a scalar for ordinary physics. The operator
(\ref{eq:hidden}) carries fermion number 2 and spin 0, and it is the order parameter of the singlet
superfluid of hypothesis H1.

\subsection{What has been established, and what has not}

Established, by exact arithmetic and by the classification of the 28 pair states: no
Lorentz-invariant pair of one 16-component field exists; a vector of pairs exists; its hidden-axis
component is a scalar of ordinary spacetime and is the unique $s$-wave pair that is a singlet of the
ordinary rotations and of the rotations of the three additional coordinates.

Not established by this section: that the hidden-axis pair is the pair that actually condenses. That
is a question of dynamics, to which Secs.~\ref{sec:vs} to \ref{sec:order} are devoted.
